\section{Introduction}
\label{sec:introduction}

Scaling inference-time computation offers a practical path to improved accuracy
without additional model training~\citep{snell2024scaling}.
By generating multiple candidate answers under a fixed budget and aggregating or
selecting among them, a system can achieve accuracy gains comparable to those from
model scaling, at inference time~\citep{wang2023self}.
In settings where the inference budget is strictly matched per method---every compared
candidate-producing method receives the same number of model calls per example---the
quality of the final answer depends entirely on the selection rule applied to the
already-generated pool.
Majority vote treats all candidates equally; but generation runs produce metadata
that is richer than raw candidate answers, and this metadata may carry signal about
which candidates are more reliable.

This paper studies that downstream decision problem explicitly: once several
candidate producers have already run under a matched-budget contract, can a final
selector use runtime-visible signals to decide when to keep one answer source and
when to defer to another? We use the term \emph{selective deferral} to emphasize
that the selector may retain the default frontier answer on most examples while
routing to an alternative source on a narrow subset identified from runtime
metadata and cross-source agreement.

We study this question in a concrete, reproducible setting.
The task is grade-school arithmetic reasoning on the GSM8K
benchmark~\citep{cobbe2021training}, the provider is Cohere command-r-plus-08-2024,
and the budget is $B = 6$ inference calls per example.
A frontier generator produces candidate answers via program-aided
execution~\citep{gao2023pal}; three independent external methods (L1, S1, and TALE)
provide reference answers under the same budget.
Given these already-generated outputs, the goal is to select a final answer without
additional inference calls and without using gold labels.

Failure-Trace Allocator (FTA) is the specific instantiated policy studied here.
FTA is a deterministic two-gate selective-deferral policy that relies on
failure-trace metadata logged during normal frontier execution plus external answer
agreement.
The \textit{Low-Depth Guard} (LDG) fires when the frontier's
\texttt{override\_reason} field signals a structurally weak single-branch output,
substituting an external majority vote.
The \textit{External-Consensus Fallback} (ECO) fires when all three external
baselines unanimously prefer a different answer than the frontier's selection.
Both gates are strictly gold-free: no correctness labels, ground-truth answers, or
example identifiers are accessed at decision time.
The two gates operate in strict priority order; at most one fires per example.

On a disjoint Aggregate-720 combining three independent evaluation seeds (seeds 41,
61, 71; $300 + 120 + 300$ examples with zero example-id or question-hash overlap),
Failure-Trace Allocator (FTA) achieves 80.69\% accuracy versus 77.64\% (L1), 77.08\% (S1), and
75.14\% (TALE).
Source-stratified 95\% bootstrap confidence intervals are strictly positive for all
four pairwise comparisons, including versus the source-local best external
(CI lower bound $+1.11$~pp).
We also report a stronger pooled-answer ensemble baseline over frontier, L1, S1,
and TALE; FTA remains ahead by point estimate on Final-300 and Aggregate-720, but
that comparison is not statistically separated.
Accordingly, the paper's claim is not broad superiority over ensembling; it is that
gold-free failure-trace metadata can identify selective recovery cases under a
matched per-method budget contract.
This contract is an evaluation design: generating frontier+L1+S1+TALE from scratch
is not equivalent to a single six-call run.
All claims are explicitly scoped to the tested provider and benchmark; we make no
universal generalization claims.

The paper's contributions are therefore organized around the general decision
problem rather than around a single implementation label:
\begin{itemize}[leftmargin=1.5em]
  \item We define a post-generation answer-selection problem under matched per-producer budgets, where candidate generation and final selection are treated as separate stages.
  \item We study a class of deterministic, gold-free selective-deferral policies that use runtime-observable metadata and cross-source agreement rather than learned post-hoc supervision.
  \item We instantiate that policy class with FTA, a two-gate routing policy using failure-trace and agreement predicates that are interpretable and runtime-legal.
  \item We present a matched-budget evaluation protocol with disjoint aggregate evidence, making explicit that the budget contract applies to candidate producers and not to constructing the full multi-source pool from scratch.
  \item We report bounded empirical evidence: in the tested Cohere/GSM8K/$B=6$ setting, FTA improves over the individual external baselines while remaining competitive with, but not statistically separated from, a strong pooled-answer baseline.
\end{itemize}

We also report five additional evaluated variants to clarify the boundary of the
supported claim and to separate exploratory development from held-out evaluation
evidence.
Their implementation labels are retained in Section~\ref{sec:ablation} and
Table~\ref{tab:negative_results} for reproducibility, but they are not part of
the final two-gate policy analyzed in the main results.

The remainder of the paper is organized as follows.
Section~\ref{sec:related_work} reviews related work on test-time compute scaling,
verifier-guided reranking, and cost-aware routing.
Section~\ref{sec:method} specifies FTA precisely, including exact trigger
conditions and a policy summary table.
Sections~\ref{sec:experiments} and~\ref{sec:results} describe the evaluation
protocol and report results with full confidence intervals.
Section~\ref{sec:ablation} summarizes ablation variants and additional evaluated outcomes.
Sections~\ref{sec:discussion}--\ref{sec:conclusion} discuss implications,
limitations, and directions for future work.
